% !TeX root = thesis.tex

\chapter{Discussion} \label{chap:discussion}

Chapter~\ref{chap:results} reported the results without interpreting them. This chapter explains what they
mean and where they are limited.

\section{Deep Learning Against the Traditional Baselines} \label{sec:disc_baselines}

The DL models beat the baselines on all three tasks. The gap is widest on PdM, where the baselines produce
a negative $R^2$ at most window sizes, meaning they do worse than always predicting the mean severity. On
AD the baselines fail mostly on recall, detecting between 22\% and 37\% of the injected anomalies against
75\% to 94\% for the DL models.

The per-fault results show where the difference comes from. Both families detect the cooling fault well,
because it drives TT-901 from around 61$^\circ$C to over 135$^\circ$C, far outside anything seen in normal
operation. On the clogged filters fault the baselines never exceed 0.294 while the DL models reach 0.84.
That fault is a slow drift in PT-903, a pressure signal that already oscillates over a comparable range
with the VPSA cycle, so no single reading and no window average looks abnormal. Only the evolution of the
sequence reveals it, and that is exactly what the summary statistics used by the baselines discard.

\section{What Temporal Context Provides} \label{sec:disc_window}

Longer windows helped less than expected. The DL models stay within a narrow band across the whole range
and already score above 0.79 on AD with a window of 1, where there is no temporal context at all. No model
reached its best result at the largest window size.

The reason is how the faults were calibrated. Each severity level is anchored to the manufacturer's
thresholds, so a fault pushes the sensor outside its normal band by construction. A single observation is
often enough to see that. The baselines behave differently, improving up to a window of 16 or 32 and then
declining, because each window is labelled by its last observation while the statistics are computed over
the whole window. The longer the window, the less those statistics describe the moment being predicted.

\section{Single-Task Against Multi-Task Learning} \label{sec:disc_stl_mtl}

No family won on all three tasks. The single-task models were best on AD and PdM, the cascaded multi-task
architecture was best on FD, and every margin was between 0.011 and 0.023. These differences are too small
to rank the approaches.

The useful finding is that one multi-task model matches three separate ones. The single-task approach needs
three models trained, tuned and maintained independently; the multi-task models share one backbone and
produce all three outputs at once. Equivalent performance at that cost favours the multi-task approach for
a real deployment. Between the two, the cascaded variant matched or beat the joint variant everywhere, with
the biggest difference on FD (0.720 against 0.640). This fits its design: its fault head is told whether an
anomaly is present before deciding which fault it is, and FD is the task that depends most on that answer.

\section{Limitations} \label{sec:disc_limitations}

No real vacuum pump fault was available. Every fault was synthetically injected from the manufacturer's
troubleshooting guide. The results show how well the models recognise the deviations that were designed,
not real degradation.

The faults were injected into only two of the eleven sessions. The two behave similarly, so the results say
nothing about generalising to another plant, another pump, or another operating regime.

Every model was evaluated at a fixed decision threshold of 0.50, rather than at a threshold chosen on
validation data. A different threshold would change the balance between precision and recall for all
classification models.

The hyperparameter search was small. Six configurations per model, with 64 hidden units as the largest
backbone, and that largest option was often selected as the best. Larger networks might do better.